% ==================================================================================================================================
% Introduction

\minitoc  % Affiche la table des matières pour ce chapitre

Depuis le lycée, nous avons définis l'intégrale avec la définition de Riemann (ou Darboux pour les plus chauvins d'netre nous).

\begin{definition}[Intégrale de Riemann]
    Soit $[a,b] \subseteq \R $ et soit $f$ une fonction continue par morceaux sur $[a,b]$. On peut alors définir :
        \[ \int_{a}^{b} f = \int_{a}^{b} f(x) dx \]
\end{definition}

Cependant cette définition se heurte à quelques problèmes dont le plus important est le fait 
qu'elle ne peut pas se généraliser à des fonctions non définies sur $[a,b]$, non continues par morceaux ou à plusieurs variables.

% ==================================================================================================================================
% Tribu Borélienne

\section{Tribu Borélienne}

\begin{definition}[Tribu]
	Soit $X$ un ensemble. Une \textbf{tribu} ou $\sigma$-algèbre sur $X$ est une partie 
	$\mathcal{B} \subseteq \mathcal{P}(X)$ telle que :
	\begin{itemize} 
		\item $ \emptyset \in \mathcal{B} $
		\item $ \forall A \in \mathcal{B}, A^c \in \mathcal{B} $
		\item Stabilité par union dénombrable : 
			$ \forall (A_i)_{i \in I}  A_i \in \mathcal{B}, \quad \forall i \in I, \quad \bigcup_{i \in I} A_i \in \mathcal{B} $
	\end{itemize}
Un élément de $\mathcal{B}$ est appelé une "partie mesurable" de $X$.
\end{definition} 

\begin{definition}[Tribu Borélienne] 
	Soit $E$ un espace topologique. La tribu borélienne est le plus petite tribu de $E$ contenant tous les ouverts de $E$. On la note $\mathcal{B}_E$. Un borélien de $E$ est une partie mesurable au sens de $\mathcal{B}_E$.
\end{definition}

\begin{example} 
	$\mathcal{B}_\R$ est la tribu borélienne, i.e la plus petite tribu contenant tous les ouverts de $\R$
\end{example} 

\begin{prop} 
	Une sigma-algèbre est stable par intersection dénombrable.
\end{prop} 

% ==================================================================================================================================
% Fonctions Mesurables

\section{Fonctions Mesurables} 

Un couple $(X,\mathcal{B})$ où $X$ est un ensemble et $\mathcal{B}$ une tribu de $X$ est appelé \textbf{espace mesurable}.

\begin{definition}[Fonction Mesurable]
	Soit $f$ une fonction. On dit que $f : (X,\mathcal{B}) \longrightarrow (X',\mathcal{B}')$ est mesurable ssi :
		\[ \boxed{ \forall A \in \mathcal{B}', \quad f^{-1}(A) \in \mathcal{B} } \]
\end{definition} 

\begin{remark}
	On appelle fonction borélienne toute fonction mesurable au sens de la tribu borélienne. 
\end{remark}

\begin{prop}
Quelques propriétés (utiles) sur la mesurabilité :
\begin{itemize} 
	\item La mesurabilité est stable par composition.
	\item La mesurabilité est stable par recollement dénombrable i.e :
		\[ f : X \longrightarrow X' \quad X = \bigcup_{n \in \N} A_n \text{ parties deux à deux disjointes mesurables} \]
		\[ \text{ si } \forall n \in \N, \quad f \restriction_A \text{ est mesurable} \]
	\item Toute fonction continue $f : E \longleftarrow E'$ ou $E$ et $E'$ sont deux espaces topologiques est mesurable.
	\item Toute fonction continue est mesurable (réciproque généralement fausse)
	\item De même, toute fonction continue par morceaux est mesurable.
\end{itemize}
\end{prop}

% ==================================================================================================================================
% Fonctions mesurables réelles ou complexes

\section{Fonctions mesurables réelles ou complexes} 

\begin{prop} 
	Soit $X$ un ensemble et $A \subseteq X$ et soit $1_A$ la fonction indicatrice de $A$.
Alors : 
		\[ \forall A \subseteq X, \quad A \in \mathcal{B} \Longleftrightarrow 1_A \in \mathcal{M}_+(X) \] 
\end{prop} 

\begin{theorem}[Stabilité]
	La mesurabilité est stable par opérations algébriques élémentaires (+,.).
\end{theorem} 

\begin{theorem}[Convergence Simple]
	Soit $(f_n)$ une suite de fonctions à valeurs dans $\mathcal{M}(X)$ qui converge simplement vers une fonction $f : X \longrightarrow \C$ 
		\[ \text{i.e } \forall x \in X, \lim_{n\to\infty} f_n(x) = f(x) \]
		% \[ \Longleftrightarrow \forall \varepsilon > 0, \forall x \in X, \exists n_0 \in \N, \forall n \geq n_0, \lvert f_n(x) - f(x) \rvert < \varepsilon \] 
Alors : $f \in \mathcal{M}(X)$ 
\end{theorem} 

\begin{theorem}[Stabilité par inf/sup]
	Soit $(f_n)_{n\in\N} \in (\overline{\mathcal{M}_+(X)})^\N$ une suite de fonctions, on a :
	\[ \sup_{n\in\N} f_n : 
		\begin{cases}
			X &\longrightarrow \overline{\R_+} \\
			x &\longmapsto \sup_{n\in\N} f_n(x) = \sup \{ f_n(x) : n \in \N \}
		\end{cases} \]
Alors :  $\sup_{n\in\N} f_n \in \overline{\mathcal{M}_+(X)}$
\end{theorem}

\begin{definition}[Fonction Etagée]
	Soit $(X,\B)$ un espace mesurable. Une fonction $e : X \rightarrow \R$ est dite étagée si :
		\begin{itemize}
			\item $e \in \M_\R(X)$
			\item $e$ prend un nombre fini de valeurs distinctes 
		\end{itemize}
\end{definition}

\begin{prop}
	$e$ est une fonction étagée ssi $e$ est une combinaison linéaire d'indicatrices de parties mesurables.
\end{prop}

\begin{theorem}[Fondamental de l'intégrale de Lebesgue]
	si $f \in \overline{\M_+(X)}$, alors il existe une suite de fonctions croissantes étagées $(e_n)_{n\in\N}$ telle que :
		\[ \boxed{ (e_n) \quad \overset{\text{CS}}{\longrightarrow} \quad f} \] 
\end{theorem}
